% Physical page 2 of the concise study, which carries this sheet and nothing
% else. The dates, relation labels, alternatives and unresolved states are the
% generated projection of research/chronology.toml, and stand exactly as the
% terminal appendix of the expansive study prints them; the location headings and
% explanatory rows are compressed for this page. The generated definitions
% are imported here, once, because the concise entrypoint's own preamble does
% not carry them.
\input{liturgy/roman-rite/postconciliar/roman-missal-third-edition-en-us-2011/propers/temporal/pc-s52-twenty-sixth-sunday-in-ordinary-time-year-a/research/chronology-annotations}

\section{Scriptural Date and Location}

\zlabel{triptych:concise:chronology:start}
\begin{passagedossiertable}
\textbf{Responsorial Psalm --- Ps 25 (24):4--5, 6--7, 8--9; response v.~6a}. No place
named in the verses\\*
\chronodate{responsorial-psalm}{\chronologyannotation{responsorial-psalm}}\\*
The Vulgate title, Douay--Rheims ``Unto the end, a psalm for David'' (24:1), attributes this alphabetic prayer for guidance and pardon without an occasion, date or placement in his reign. The NABRE bound is the Psalter's; no individual psalm is securely dated.\\
\midrule
\textbf{Communion Antiphon, first option --- cf.\ Ps 119 (118):49--50, an adaptation}. No place named in the verses\\*
\chronodate{communion-antiphon-a}{\chronologyannotation{communion-antiphon-a}}\\*
The \latin{Zain} strophe belongs to the alphabetic psalm on God's law. Its Vulgate title, \latin{Alleluia}, names no author; Haydock reports general Davidic attribution and competing occasions. The reign identifies David; composition has the common critical bound.\\
\midrule
\textbf{First Reading --- Ezek 18:25--28}. Among the captives ``by the river Chobar''
(1:1--3); addressed to ``the house of Israel'' (18:25)\\*
\chronodate{first-reading}{\chronologyannotation{first-reading}}\\*
Ezekiel is the priest, son of Buzi (1:3). The undated chapter stands between the sixth- and seventh-year oracles (8:1; 20:1). The setting is Joachin's captivity: the preferred label comes from 4~Kgs 24:12; Haydock reports Usher's A.M. figure; Reid dates surrender and Schets deportation a year earlier (\work{Catholic Encyclopedia}). Schets reads the prophecy's duration from 29:17. None dates this oracle.\\
\midrule
\textbf{Entrance Antiphon --- Dan 3:29--31, 42--43 (Vulgate), adapted}. Nabuchodonosor's
furnace, ``of the province of Babylon'' (3:1, 20--25)\\*
\chronodate{entrance-antiphon}{\chronologyannotation{entrance-antiphon}}\\*
\textit{Narrated event: the prayer of Azarias in the midst of the fire.}\\*
The Missal orders these verses 31, 29, 30, 43, 42, from the Greek/Vulgate addition absent from Hebrew Daniel. The setting is Jeremiah's foretold service to Babylon (Jer 25:11--12), not a year for this scene. Gigot (\work{Catholic Encyclopedia}, ``Daniel'') ascribes 3:24--90 to an unknown Jewish author and reports that nearly all Catholic writers of his day placed the prayer in Daniel's time.\\
\midrule
\textbf{Gospel --- Mt 21:28--32}. In the temple at Jerusalem, to ``the chief priests and
ancients of the people'' (21:23)\\*
\chronodate{gospel}{\chronologyannotation{gospel}}\\*
\textit{Narrated event: a parable spoken after the question about Jesus'
authority and John's baptism (21:23--27).}\\*
Traditional attribution: Matthew the Apostle; the event is undated here. Jacquier's disputed composition dates (\work{Catholic Encyclopedia}) report ancient writers and divided Catholic critics; Durand's last alternative dates an Aramaic original. These dates concern the book.\\
\midrule
\textbf{Alleluia verse --- Jn 10:27, with a speaker formula added}. The temple, ``in
Solomon's porch'', at the Dedication, in winter (10:22--23)\\*
\chronodate{gospel-acclamation}{\chronologyannotation{gospel-acclamation}}\\*
\textit{Narrated event: Jesus' answer to those who asked him to say
plainly whether he was the Christ.}\\*
Traditional attribution: John the Apostle. Durand dates the Johannine writings; Fonck reports the general opinion while conceding uncertain historical information (\work{Catholic Encyclopedia}). These composition ranges do not date the discourse.\\
\midrule
\textbf{Second Reading, longer form --- Phil 2:1--11}. From Paul ``in my bands''
(1:7, 13), at Rome in the sources cited; to the saints ``at Philippi'' (1:1)\\*
\chronodate{second-reading-long}{\chronologyannotation{second-reading-long}}\\*
\textbf{Second Reading, shorter form --- Phil 2:1--5}. As for the longer form\\*
\chronodate{second-reading-short}{\chronologyannotation{second-reading-short}}\\*
Traditional attribution: Paul in captivity. Prat dates the captivity letters; Vander Heeren gives a Roman range and reports disagreement over an early or late point in that stay (\work{Catholic Encyclopedia}).\\
\midrule
\textbf{Communion Antiphon, second option --- 1~Jn 3:16}. No place named in the letter\\*
\chronodate{communion-antiphon-b}{\chronologyannotation{communion-antiphon-b}}\\*
Traditional attribution: John the Apostle. Durand dates the Johannine writings; Drum finds Ephesus and the first century's last years probable (\work{Catholic Encyclopedia}).\\
\end{passagedossiertable}
\zlabel{triptych:concise:chronology:end}
